\section{Опис виконання роботи}

\subsection{Загальна структура}

Система \projname{} має чотирирівневу архітектуру (рис.~\ref{fig:arch}):

\begin{verbatim}
+----------------------------------------------------------+
| Operator tier                                            |
| missions/relay_demo.py  missions/inspection.py           |
| swarm/operator.py  (CLI: takeoff / goto / scan / ...)    |
+------------------------+---------------------------------+
                         | JSON commands (MQTT v5, QoS 1)
                         v
+----------------------------------------------------------+
| Drone tier (Isaac Sim + Pegasus  OR  mock_swarm)         |
| swarm/controller.py  (CommandedNonlinearController)      |
|   +-- swarm/trajectory.py  (quintic segments)            |
|   +-- swarm/coverage.py    (boustrophedon path)          |
|   +-- swarm/relay.py       (DroneStateTracker, Relay)    |
|   +-- Pegasus NonlinearController (SE(3) geometric ctrl) |
+------------------------+---------------------------------+
                         | state + telemetry (MQTT, QoS 0)
                         v
+----------------------------------------------------------+
| Cloud tier                                               |
| Mosquitto  -->  Telegraf  -->  InfluxDB v2               |
+------------------------+---------------------------------+
                         | InfluxDB Flux query
                         v
+----------------------------------------------------------+
| Visualization tier                                       |
| Grafana dashboard (operator workstation)                 |
+----------------------------------------------------------+
\end{verbatim}

\subsection{Модулі системи}

\begin{table}[h]
  \centering
  \caption{Модулі програмного забезпечення}
  \begin{tabularx}{\linewidth}{llX}
    \toprule
    Файл & Модуль & Призначення \\
    \midrule
    \texttt{swarm/commands.py}       & Commands   & Pydantic-схема 7 типів команд; валідація payload \\
    \texttt{swarm/config.py}         & Config     & \texttt{MqttSettings}, \texttt{SimSettings} — читаються з \texttt{.env} \\
    \texttt{swarm/facility.py}       & Facility   & POI, ScanZone, \texttt{InspectionMission} \\
    \texttt{swarm/trajectory.py}     & Trajectory & QuinticSegment, TrajectoryFollower, build\_orbit \\
    \texttt{swarm/coverage.py}       & Coverage   & \texttt{boustrophedon\_path} — BCD-waypoints \\
    \texttt{swarm/controller.py}     & Controller & \texttt{CommandedNonlinearController} (підклас Pegasus) \\
    \texttt{swarm/relay.py}          & Relay      & \texttt{DroneStateTracker}, \texttt{RelayMission} \\
    \texttt{swarm/mock\_swarm.py}    & Mock       & Headless-симулятор для тестування без Isaac Sim \\
    \texttt{swarm/mqtt\_publisher.py}& MQTT Pub   & Публікація команд, стану, телеметрії, подій \\
    \texttt{isaac/swarm\_app.py}     & Isaac App  & Ініціалізація Pegasus, спавн дронів, камери \\
    \texttt{missions/relay\_demo.py} & Mission    & Демо естафетної передачі двома дронами \\
    \texttt{missions/inspection.py}  & Mission    & Повна інспекція трьома дронами \\
    \texttt{missions/fly\_demo.py}   & Mission    & Демо прямолінійного польоту одним дроном \\
    \bottomrule
  \end{tabularx}
  \label{tab:modules}
\end{table}

\subsection{Стек керування на основі літератури}

Ієрархія \emph{operator – planner – tracker – plant}:

\begin{verbatim}
Operator CLI (swarm.operator)
        |  JSON command  -- MQTT v5 [17] -->
        v
swarm.commands (Pydantic)
        v
swarm.controller.CommandedNonlinearController
        |
        +-- swarm.trajectory.TrajectoryFollower
        |     +-- single quintic segment      [Lynch&Park 2017]
        |     |       (takeoff, goto, hover, land)
        |     +-- chained quintic segments    [Lynch&Park 2017]
        |             through Boustrophedon waypoints [Choset 2000]
        |
        +-- Pegasus NonlinearController (без модифікацій)
                +-- outer-loop position PD+I           -+
                +-- SO(3) attitude PD                    +- [Mellinger&Kumar 2011]
                +-- force/torque -> rotor allocation     |  [Pinto et al. 2021]
                                                         |  (анкер: [Lee 2010])
                                                         v
                                                Pegasus Multirotor (PhysX)
\end{verbatim}

\subsection{MQTT-топіки}

\begin{table}[h]
  \centering
  \caption{Топіки MQTT v5}
  \begin{tabularx}{\linewidth}{lllX}
    \toprule
    Топік & Напрям & QoS & Зміст \\
    \midrule
    \texttt{swarm/cmd/\{id\}} & Op $\to$ Drone & 1 & JSON-команда \\
    \texttt{swarm/cmd/all}    & Op $\to$ All   & 1 & Broadcast \\
    \texttt{swarm/state/\{id\}}     & Drone $\to$ Cloud & 0 & Позиція, mode, battery\_pct \\
    \texttt{swarm/telemetry/\{id\}} & Drone $\to$ Cloud & 0 & InfluxDB Line Protocol \\
    \texttt{swarm/events/\{kind\}}  & Drone $\to$ Cloud & 1 & command\_accepted, trajectory\_finished \\
    \bottomrule
  \end{tabularx}
  \label{tab:topics}
\end{table}

QoS~1 (at-least-once delivery з PUBACK-handshake) використовується для команд та подій; QoS~0 для телеметрії, де мінімальна затримка важливіша за гарантію доставки~\cite{MQTT5}.

\subsection{Машина станів дрона}

Кожен дрон реалізує просту машину станів:

\begin{center}
\texttt{idle} $\xrightarrow{\text{takeoff}}$ \texttt{takeoff} $\xrightarrow{\text{finish}}$ \texttt{hover}
$\xrightarrow{\text{goto/scan/orbit}}$ \texttt{(mode)} $\xrightarrow{\text{finish}}$ \texttt{hover}
$\xrightarrow{\text{land}}$ \texttt{land} $\xrightarrow{\text{finish}}$ \texttt{idle}
\end{center}

Команда \texttt{stop} на будь-якому кроці переводить дрон у \texttt{hover} на поточній позиції.
Поточний \texttt{mode} публікується в \texttt{swarm/state/\{id\}} кожні 0,2~с.
